% Chapter 3 worked example. Included by ch03_resources.tex.
\section{Worked Example: Hospital Command Dashboard and Patient Flow}
\label{sec:comp-ch03-worked-example}

\subsection{Purpose and Status}
This worked example demonstrates how the methods in Chapter~3 can be
translated into a reviewable institutional decision. All numerical inputs are
synthetic unless a source is identified explicitly. They are intended to show the
logic of the analysis, not to report empirical findings or establish a universal
threshold. Local use requires replacement of the assumptions with current,
versioned evidence and approval through the relevant clinical, managerial,
economic, legal, technical, and governance processes.

A tertiary hospital plans a command dashboard combining admissions, discharges, bed status, staffing, diagnostics, transport, and forecast demand. Leaders expect faster decisions, but current indicators have inconsistent definitions and no explicit action owners. 

\subsection{Decision Problem and Comparators}
\textbf{Decision problem.} Should the hospital implement an AI-enhanced command dashboard, redesign the existing operational review process, or combine both through a staged signal-to-action intervention?

The comparator set is deliberately broader than the existing pathway. This prevents
a new technology from receiving a lower evidentiary threshold than staffing,
workflow, shared-service, conventional digital, or no-change alternatives.

\begin{longtable}{@{}p{0.08\textwidth}p{0.84\textwidth}@{}}
\caption{Comparators for the Chapter 3 worked example}
\label{tab:comp-ch03-comparators}\\
\toprule
\textbf{Option} & \textbf{Comparator or strategy} \\
\midrule
\endfirsthead
\toprule
\textbf{Option} & \textbf{Comparator or strategy} \\
\midrule
\endhead
\bottomrule
\endlastfoot
1 & Current morning operational meeting and static reports \\
2 & Standardized dashboard without predictive functions \\
3 & AI-enhanced dashboard with demand forecasts and anomaly detection \\
4 & Targeted process redesign for discharge, imaging, and transport bottlenecks \\
\end{longtable}

\subsection{Model and Decision Structure}
The analytical structure should follow the decision rather than the available
software. The team should map the causal pathway from intervention input to
information, human decision, action, outcome, resource consequence, and review.
The structure should also identify capacity constraints, uncertainty, subgroup
variation, implementation requirements, and events that would change the
recommendation.

A compact quantitative relationship used in this example is:
\begin{equation}
T_{\mathrm{action}}=T_{\mathrm{delivery}}+T_{\mathrm{review}}+T_{\mathrm{response}}
\label{eq:comp-ch03-key-relationship}
\end{equation}
The equation is an organizing aid rather than a substitute for disaggregated evidence,
qualitative findings, or mandatory safeguards. Its variables and interpretation must
be defined in the local protocol.

\subsection{Synthetic Parameters and Evidence Sources}
Table~\ref{tab:comp-ch03-parameters} provides the base-case inputs. A
production analysis should add parameter distributions, correlations, structural
alternatives, data dates, system versions, and evidence-confidence ratings.

\begin{longtable}{@{}p{0.32\textwidth}p{0.22\textwidth}p{0.38\textwidth}@{}}
\caption{Synthetic parameters for the Chapter 3 worked example}
\label{tab:comp-ch03-parameters}\\
\toprule
\textbf{Parameter} & \textbf{Base value} & \textbf{Source or interpretation} \\
\midrule
\endfirsthead
\toprule
\textbf{Parameter} & \textbf{Base value} & \textbf{Source or interpretation} \\
\midrule
\endhead
\bottomrule
\endlastfoot
Average daily admissions & 312 & Operations data \\
Mean bed occupancy & 94\% & Capacity dashboard \\
Median data latency & 95 minutes & Interface audit \\
Indicators without named action owner & 42\% & Dashboard inventory \\
Forecast mean absolute percentage error & 13\% & Retrospective validation \\
Estimated annual dashboard cost & EUR 265,000 & Lifecycle-cost estimate \\
\end{longtable}

\subsection{Step-by-Step Analysis}
The sequence below is intended to make the analysis reproducible and to ensure
that evidence, economics, governance, and implementation develop together.

\begin{longtable}{@{}p{0.07\textwidth}p{0.55\textwidth}p{0.30\textwidth}@{}}
\caption{Analytical sequence}
\label{tab:comp-ch03-analysis-steps}\\
\toprule
\textbf{Step} & \textbf{Required work} & \textbf{Decision record} \\
\midrule
\endfirsthead
\toprule
\textbf{Step} & \textbf{Required work} & \textbf{Decision record} \\
\midrule
\endhead
\bottomrule
\endlastfoot
1 & Define the operational decisions, decision windows, and accountable owners before selecting indicators. & Evidence and responsible owner to be recorded \\
2 & Separate observed values, model estimates, forecasts, and recommendations. & Evidence and responsible owner to be recorded \\
3 & Create an indicator dictionary with numerator, denominator, exclusions, latency, and limitations. & Evidence and responsible owner to be recorded \\
4 & Map each signal to a feasible action and at least one balancing measure. & Evidence and responsible owner to be recorded \\
5 & Test forecast value against a standardized non-predictive dashboard and targeted process redesign. & Evidence and responsible owner to be recorded \\
6 & Measure action latency, downstream congestion, workload, and the consequences of recommendations that cannot be implemented. & Evidence and responsible owner to be recorded \\
\end{longtable}

\subsection{Illustrative Outputs}
The outputs in Table~\ref{tab:comp-ch03-outputs} show how technical,
clinical, operational, economic, equity, and governance findings should be kept
visible rather than compressed into a single favourable or unfavourable score.

\begin{longtable}{@{}p{0.30\textwidth}p{0.24\textwidth}p{0.38\textwidth}@{}}
\caption{Illustrative model and decision outputs}
\label{tab:comp-ch03-outputs}\\
\toprule
\textbf{Output} & \textbf{Result} & \textbf{Interpretation} \\
\midrule
\endfirsthead
\toprule
\textbf{Output} & \textbf{Result} & \textbf{Interpretation} \\
\midrule
\endhead
\bottomrule
\endlastfoot
Principal failure in current design & No action ownership & Information does not reliably change decisions \\
Forecast technical performance & Potentially useful & Requires local operational validation \\
Main value constraint & High occupancy and limited downstream capacity & Forecasting cannot create beds or staff \\
Preferred intervention & Signal-to-action redesign plus staged dashboard & Dashboard alone is insufficient \\
Primary monitoring outcome & Decision and action latency & Views and logins are secondary \\
\end{longtable}

\subsection{Sensitivity, Uncertainty, and Subgroups}
At minimum, the completed analysis should vary the parameters most likely to
change the decision, test at least one credible alternative model structure, and
report subgroup results separately where performance, access, response, burden,
or opportunity cost may differ. Deterministic analysis should identify thresholds;
probabilistic analysis should be used where credible parameter distributions and
correlations are available. Scenario analysis is preferable when structural or
future uncertainty cannot be represented defensibly by one probability model.

The record should distinguish uncertainty that can be reduced through evidence,
uncertainty that can be managed through safeguards, and uncertainty that remains
too consequential to accept. Missing evidence should not be represented as an
average or neutral value without justification.

\subsection{Interpretation and Recommendation}
Implement a standardized operational dashboard first, with explicit action ownership and decision thresholds. Add forecasting functions only after the data-latency and response-pathway requirements have been met and tested in simulation and a controlled pilot.

The recommendation is conditional rather than permanent. It applies only to the
specified population, setting, intervention configuration, evidence cut-off, and
version. The following conditions should be incorporated into the approval,
contract, implementation, monitoring, or renewal record:

\begin{itemize}
 \item all priority indicators have definitions, owners, decision windows, and balancing measures.
 \item data latency remains within the operational window.
 \item forecast uncertainty is displayed and observed values are distinguishable from predictions.
 \item the pilot measures action and outcome changes rather than views or user satisfaction alone.
 \item a governance process reviews gaming, surveillance concerns, and unequal workload effects.
\end{itemize}

\subsection{Decision Statement}
\begin{quote}
For the specified decision problem and comparator set, the institution should
\emph{[proceed / proceed with conditions / redesign / pilot / restrict / defer /
replace / retire]}. The decision applies to \emph{[population, setting, version,
and evidence period]}, is owned by \emph{[accountable authority]}, and will be
reviewed on \emph{[date]} or earlier if \emph{[trigger events]} occur. The
evidence, assumptions, unresolved uncertainty, mandatory safeguards, monitoring
thresholds, and exit arrangements are recorded in the accompanying dossier.
\end{quote}
